% Alphabet invariance: paper-form write-up of the transfer theorem.
% Self-contained section; labels are fresh (def:code, lem:dict, lem:selfsync,
% lem:pass-transfer, def:conjugate, thm:transfer-forward, lem:compose,
% lem:retract, thm:transfer, cor:uniform, lem:palindrome, thm:rev-uniform,
% rem:native, prop:unary-edge). NOT integrated into main.tex; the constants
% and conventions follow the paper's Sections 2-4.

\section{Alphabet Invariance}
\label{sec:alphabet}

Every construction so far has fixed its alphabet, and the two letters always
sufficed. This section addresses the fourth open problem of
Section~\ref{sec:calculus} head on: \emph{which functions are reachable over
every alphabet with $|\Sigma| \ge 2$, and does a larger alphabet ever add
power?} The answer is that, up to a canonical conjugation, the reachable
class does not depend on the alphabet at all: for any two alphabets of at
least two letters and any \emph{good coding} between them, $f$ is reachable
over the one iff its conjugate is reachable over the other. The boundary of
the invariance is exactly $|\Sigma| = 1$: the conjugation exists one way and
cannot exist the other way, for a structural reason.

\subsection{Comma-Free Codings}

\begin{definition}
    (\emph{Good Codings})
    \label{def:code}
    A \emph{coding} is an injective monoid homomorphism $c : \Sigma^*
    \rightarrow \Gamma^*$ with $D_c \triangleq c(\Sigma)$ of uniform length
    $\ell \ge 1$ (a \emph{dictionary}); write $c(S)$ for the letter-by-letter
    extension. The coding is \emph{good} if $D_c$ is \emph{comma-free}
    (Golomb's condition; the theory of codes treats it under synchronizing
    codes \cite{berstel85,lothaire97}): no $w \in D_c$ occurs in $u \cdot v$
    at any of the junction-straddling positions $1, \dots, \ell-1$
    ($u, v \in D_c$).
\end{definition}

Comma-freeness is the Golomb analogue of the Comma Code Lemma for
$\mathtt{enc}^2_{b,x}$ (blocks $xc$): it is exactly the self-synchronization
that makes the coding usable inside a scanning calculus. Without it the
conjugacy class of a function depends on the coding, and the notion
degenerates: a free assignment $\ell$ with $\ell(SS) = \ell(S)^2$, say,
makes $X \mapsto XX$ conjugate to the unary squaring map. The good codings
are the ones on which the conjugation is canonical.

\begin{lemma}
    (\emph{Existence: the Double-$a$ Family})
    \label{lem:dict}
    Let $\Gamma$ have two distinct letters $a \ne b$ and $\ell \ge 3$. Then
    \begin{align}
        D_\ell \;\triangleq\; \{\, a\,a\,m\,b \;:\; m \in
        \Gamma^{\ell-3},\ aa \not\le m,\ m[0] \ne a \,\}
    \end{align}
    is comma-free, of uniform length $\ell$; over the binary alphabet
    $|D_\ell| = \mathrm{Fib}(\ell-2)$, so good codings exist between any two
    alphabets with at least two letters.
\end{lemma}

\begin{proof}
    Every codeword starts with $aa$, contains $aa$ only at position $0$
    (the interior $m$ avoids $aa$ and, for $\ell > 3$, does not start with
    $a$), and ends with $b$. In $u \cdot v$ the factor $aa$ therefore occurs
    only at $0$ and $\ell$ (the junction reads $b \cdot a$), and any
    occurrence of a codeword, starting as it does with $aa$, starts at $0$
    or $\ell$ -- aligned. The count: $m$ is $b$ followed by an $aa$-free
    word of length $\ell-4$, counted by $\mathrm{Fib}(\ell-2)$.
    (Verified: comma-freeness of all 21 dictionaries in the family, binary/
    ternary/ quaternary, $\ell = 3..9$; sizes at $\ell = 3..12$ match
    $\mathrm{Fib}(\ell-2)$ with $\mathrm{Fib}(1) = \mathrm{Fib}(2) = 1$.)
\end{proof}

\begin{lemma}
    (\emph{Self-Synchronization})
    \label{lem:selfsync}
    Let $c$ be good and $B \ne \epsilon$. The occurrences of $c(B)$ in
    $c(C)$ are exactly $\{\, \ell \cdot s \;:\; s$ an occurrence of $B$ in
    $C \,\}$; in particular every occurrence of a codeword inside a
    codeword-sequence starts at a multiple of $\ell$.
\end{lemma}

\begin{proof}
    Uniform lengths first: an occurrence of $w \in D_c$ (length $\ell$)
    lying wholly inside a single codeword (length $\ell$) of the surrounding
    sequence starts at offset $0$. Otherwise the occurrence straddles at
    least one junction $j\ell$ with $p < j\ell < p + \ell$, and its
    restrictions to the two flanking codewords exhibit $w$ in their product
    at the junction-straddling position $p - (j{-}1)\ell \in [1, \ell-1]$ --
    against comma-freeness. For the block $c(B)$: its first codeword $c(B[0])
    \in D_c$ occurs at a multiple of $\ell$ by the first part, the block
    then covers whole codewords of $c(C)$, and these spell $B$; the converse
    is immediate. (Verified: all $B \in \{a,b,c\}^{\le 3} \setminus
    \{\epsilon\}$, $C \in \{a,b,c\}^{\le 5}$ -- 14,196 cases, zero
    exceptions.)
\end{proof}

\begin{lemma}
    (\emph{Pass Transfer})
    \label{lem:pass-transfer}
    Let $c$ be good. For all $A, C \in \Sigma^*$ and $B \in \Sigma^*
    \setminus \{\epsilon\}$,
    \begin{align}
        [\,c(A)/c(B)\,]\; c(C) \;=\; c\big([\,A/B\,]C\big),
    \end{align}
    and the same identity holds for the once, rightmost, and $k$-th
    occurrence primitives: the greedy occurrence lists correspond
    order-isomorphically under $s \mapsto \ell s$.
\end{lemma}

\begin{proof}
    Induction along the scan. By Self-Synchronization the occurrence lists
    correspond, and the correspondence preserves order and disjointness
    (positions scale by $\ell$), so leftmost, $k$-th and rightmost all
    transfer; it is direction-free, so the right-to-left variant transfers
    too. For the full pass, write $C = X B Y$ with $|X|$ minimal: the first
    occurrence of $c(B)$ in $c(X)\,c(B)\,c(Y)$ is the aligned one at
    $\ell|X|$ (any earlier one would be an aligned occurrence inside
    $c(X)$, i.e.\ an occurrence of $B$ in $X$), the scan resumes at
    $\ell(|X|{+}|B|)$ -- aligned again, and the inserted $c(A)$ is never
    re-entered, matching the never-rescan clause of Definition~\ref{def:subst}.
    Both sides now split as $c([A/B]X) \cdot c(A) \cdot c([A/B]Y)$ with
    $|X|, |Y| < |C|$. (Verified: $A, C \in \{a,b,c\}^{\le 3}$, $B \in
    \{a,b,c\}^{1..3}$ -- 188,760 leftmost cases; once, rightmost and
    $k$-th ($k \le 3$) on the smaller box -- 31,200 cases; zero mismatches.)
\end{proof}

\subsection{The Transfer Theorem}

\begin{definition}
    (\emph{Conjugate})
    \label{def:conjugate}
    For $f$ an $n$-ary partial function on $\Sigma^*$ and $c : \Sigma
    \rightarrow \Gamma^\ell$ a good coding, the \emph{conjugate}
    $f^c : \mathrm{image}(c)^n \rightharpoonup \mathrm{image}(c)$ is
    $c \circ f \circ c^{-1}$ -- a partial function in the sense of the
    restriction clause of Definition~\ref{def:reachable}.
\end{definition}

\begin{theorem}
    (\emph{Direction 1})
    \label{thm:transfer-forward}
    Let $E \in \mathrm{Exp}_n$ over $\Sigma$ and $c$ good. Define $E^c$ by:
    constants $W \mapsto c(W)$, variables and both constructors unchanged.
    Then
    \begin{align}
        \llbracket E^c \rrbracket\big(c(\vec S)\big) \;=\;
        c\big(\llbracket E \rrbracket(\vec S)\big)
    \end{align}
    whenever the right-hand side is defined, and the two sides are
    undefined together (the pattern is $\epsilon$ on one side iff on the
    other). Moreover $|E^c| = |E|$ (node count; constant symbols grow
    $\ell$-fold in length), $\deg(E^c) = \deg(E)$, and
    $\mathrm{Safe}(E) \Rightarrow \mathrm{Safe}(E^c)$; and if $\llbracket E
    \rrbracket$ is total then a \emph{total} $\Gamma$-expression agreeing
    with $E^c$ on the image exists -- wrap the pattern once:
    $[R^c/\gamma(P^c)]\,E^c$ with $\gamma(X) \triangleq
    \mathtt{if}(\mathtt{eq}(X,\epsilon), b, X)$.
\end{theorem}

\begin{proof}
    Induction on $E$. Variables and constants are immediate; concatenation
    commutes with $c$ because $c$ is a homomorphism. For the pass node the
    patterns and replacements are the $c$-images of the sub-values, and
    Pass Transfer applies at every value of the sub-expressions; the
    definedness clause is the $\epsilon$-condition, and $c(B) = \epsilon$
    iff $B = \epsilon$. Safety and degree are structural ($c(W)$ is a
    nonempty constant iff $W$ is; $\deg$ reads the replacement, and the
    replacement's image has the same degree). (Verified: all 2,954
    depth-$\le 1$ expressions over $\{a,b,c\}$ with constants $\le 2$, on
    all 121 inputs of length $\le 4$ -- 357,434 evaluations, 0 mismatches,
    23,912 undefined on both sides; 48,000 further evaluations of random
    depth-$\le 3$ expressions; the same under once and rightmost semantics,
    118,160 evaluations each; degree and safety over the full space, 2,954
    expressions, 0 violations; node count equal in every case.)
\end{proof}

\begin{lemma}
    (\emph{Composition})
    \label{lem:compose}
    If $c : \Sigma \rightarrow \Gamma^{\ell_c}$ and $d : \Gamma \rightarrow
    \Sigma^{\ell_d}$ are good, then $d \circ c$ is good: $\{d(w) : w \in
    D_c\}$ is comma-free over $\Sigma$, of uniform length
    $\ell_c\ell_d$.
\end{lemma}

\begin{proof}
    An occurrence of $d(w)$ inside a sequence of $d$-images of
    $D_c$-codewords begins with $d(w[0]) \in D_d$, which by
    Self-Synchronization over $D_d$ starts on the $D_d$-grid; the whole
    block $d(w)$ then covers $D_d$-codewords spelling
    $d(w[0]) \cdots d(w[\ell_c{-}1])$, so by injectivity of $d$ the word $w$
    occurs among the pre-image codewords, and if the block did not start on
    the outer ($\ell_c\ell_d$-) grid, $w$ would occur in a product of two
    $D_c$-codewords at a junction-straddling position -- against the
    comma-freeness of $D_c$. (Verified: $\{d(c(\sigma))\}$ for $c :
    \{a,b,c\} \rightarrow \{x,y\}^6$ and $d : \{x,y\} \rightarrow
    \{a,b,c\}^4$, length-24 dictionary, comma-free.)
\end{proof}

\begin{lemma}
    (\emph{Retraction})
    \label{lem:retract}
    Let $h = d \circ c : \Sigma \rightarrow \Sigma^{\ell_c\ell_d}$ with
    $c, d$ good. Then $h$ is good, and both $h$ and a decoder $H_h$ with
    $H_h \circ h = \mathrm{id}$ are \emph{total} $L$-reachable: $h$ is
    $\mathrm{rep}_{|\Sigma|}$ with the single characters as patterns and
    the $h$-codewords as replacements, $H_h$ the same with the roles
    exchanged -- $\mathrm{rep}_n$ places no restriction on its patterns
    (Theorem~\ref{lem:multiple substitution}).
\end{lemma}

\begin{proof}
    Comma-freeness is Lemma~\ref{lem:compose}; on an $h$-image the codeword
    occurrences are aligned (Self-Synchronization), so the freezing
    semantics of $\mathrm{rep}_n$ tiles the text exactly, replacing each
    $h(\sigma)$ by $\sigma$. Both $\mathrm{rep}_n$ have constant, nonempty
    patterns, hence are total. (Verified: the $h$- and $H_h$-expressions
    against their intended semantics on all 40 inputs of length $\le 3$
    over $\{a,b,c\}$; and as part of the roundtrip below.)
\end{proof}

\begin{theorem}
    (\emph{Transfer})
    \label{thm:transfer}
    For all finite $\Sigma, \Gamma$ with $|\Sigma|, |\Gamma| \ge 2$, every
    good coding $c$, and every partial function $f$ on $\Sigma^*$:
    \begin{align}
        f \text{ is $L_\Sigma$-reachable} \;\iff\; f^c \text{ is
        $L_\Gamma$-reachable},
    \end{align}
    and when $f$ is total the $\Gamma$-witness may be taken total. The
    $\Gamma$-witness of Direction 2 is $H_h\big(E'_d(h(X_1), \dots,
    h(X_n))\big)$ with $E' = E^c$, $E'_d = (E^c)^d$, $h = d \circ c$; it has
    $\deg$ equal to $\deg(E)$ and node count $|E|$ plus a constant depending
    only on the two alphabets.
\end{theorem}

\begin{proof}
    Direction 1 is Theorem~\ref{thm:transfer-forward}. For Direction 2 let
    $E'$ over $\Gamma$ witness $f^c$. Every $h(X_i) = d(c(X_i))$ lies in
    $\mathrm{image}(d)$, where $E'_d$ computes $d$ of the $E'$-value
    (Direction 1, for the coding $d$); so
    \begin{align*}
        \llbracket E'_d \rrbracket\big(h(\vec X)\big)
        \;=\; d\big(\llbracket E' \rrbracket(c(\vec X))\big)
        \;=\; d\big(c(f(\vec X))\big) \;=\; h\big(f(\vec X)\big),
    \end{align*}
    and $H_h$ returns $f(\vec X)$ by the Retraction lemma. Definedness
    tracks $f$'s: $h$ and $H_h$ are total, and $E'_d$ is defined at the
    $h$-images exactly when $E$ is at $\vec X$. The degree and node claims
    follow from Theorem~\ref{thm:transfer-forward} and the constant patterns
    of the two $\mathrm{rep}_n$. (Verified: the full roundtrip $E_0 \mapsto
    E' \mapsto E_{\mathrm{final}}$ on all 155 depth-$\le 1$ expressions
    with constants $\le 1$, on all 40 inputs of length $\le 3$ -- 6,200
    evaluations, 0 mismatches, 1,025 undefined on both sides; $\deg
    E_{\mathrm{final}} = \deg E_0$ in all 155 cases, node overhead 528 for
    the $(3,2)$-alphabet pair.)
\end{proof}

\begin{corollary}
    (\emph{Alphabet-Uniform Questions})
    \label{cor:uniform}
    The right pass $\mathrm{r}$, the once-hinge $\mathrm{r}_1$, and every
    $k$-th-occurrence node are \emph{coding-equivariant} -- they commute
    with every good coding, by Lemma~\ref{lem:pass-transfer} -- so each of
    their reachability questions is a single question, not one per
    alphabet: reachable over one alphabet with at least two letters iff
    reachable over every. The same holds for $\mathrm{rev}$ by
    Theorem~\ref{thm:rev-uniform} below. Consequently every negative
    search of Sections~\ref{sec:calculus}--\ref{sec:variants} run over one
    alphabet excludes the conjugate over \emph{every} alphabet of at least
    two letters.
\end{corollary}

\begin{proof}
    Equivariance plus Theorem~\ref{thm:transfer}: if $F$ commutes with the
    good coding $e$, then a $\Sigma$-witness for $F$ conjugates along $c$
    to a $\Gamma$-function which, composed with the retraction $q = c
    \circ e$, computes the native $F$ over $\Gamma$ -- the identical
    construction as in Theorem~\ref{thm:rev-uniform}, with the equivariance
    of Lemma~\ref{lem:pass-transfer} in place of the palindrome clause.
    (Verified end to end: the pass node over $\{a,b,c\}$ transferred to
    $\{x,y\}$ and retracted -- a full binary expression in the leftmost
    calculus, 294 evaluations, 0 mismatches; once, rightmost and second
    occurrence as semantic compositions, 882 evaluations, 0 mismatches;
    equivariance of all four passes under the retraction codings, 2,352
    cases, 0 mismatches; the doubling function as a further end-to-end
    witness, 31 inputs.)
\end{proof}

\subsection{Reversal}

Reversal is the one equivariance that codings do not enjoy for free: a
letter-by-letter reversal reverses \emph{each codeword individually}, so
$c(\mathrm{rev}\,S) = \mathrm{rev}(c(S))$ holds iff every codeword is a
palindrome -- a false claim in an earlier draft of this section, caught by
the verification battery (60 failures on 31 inputs). The naive ``rev on the
code image'' is not the reversal of the image; the repair is to retract
along a coding whose outer half is palindromic.

\begin{lemma}
    (\emph{Palindromic Dictionaries})
    \label{lem:palindrome}
    \begin{enumerate}[(i)]
        \item $c(\mathrm{rev}\,S) = \mathrm{rev}(c(S))$ for all $S$ iff
        every codeword of $c$ is a palindrome; then
        $e(\mathrm{rev}\,T) = \mathrm{rev}(e(T))$ for every such $e$.
        \item (\emph{Middle-marker family.}) For $\mu \in \Sigma$ and
        $m \ge 1$, $\{\, \alpha\,\mu\,\mathrm{rev}(\alpha) : \alpha \in
        (\Sigma \setminus \mu)^m \,\}$ is comma-free and palindromic, of
        size $(|\Sigma|-1)^m$ -- unbounded for $|\Sigma| \ge 3$.
        \item Over the binary alphabet the family degenerates (size 1),
        but size-3 palindromic comma-free dictionaries exist:
        $\{\mathtt{aabaa}, \mathtt{ababa}, \mathtt{abbba}\}$ and
        $\{\mathtt{aabaa}, \mathtt{babab}, \mathtt{bbabb}\}$.
    \end{enumerate}
\end{lemma}

\begin{proof}
    (i) is immediate on one-letter words. (ii) $\mu$ occurs in each codeword
    only at the exact middle, and in $u \cdot v$ only at the two middles;
    an occurrence of a codeword has its $\mu$ at $p + m$, hence $p \in
    \{0, 2m+1\}$ -- aligned. (iii) by inspection. (Verified: (ii) 20
    dictionaries, ternary and quaternary, $m \le 5$, all comma-free with the
    stated sizes; (iii) both dictionaries comma-free; exact maxima of
    palindromic comma-free sets over binary, by branch and bound:
    $1, 1, 3, 2, 7, 14$ at $\ell = 3, 4, 5, 6, 7, 9$ -- data, not a
    formula.)
\end{proof}

\begin{theorem}
    (\emph{Reversal Uniformity})
    \label{thm:rev-uniform}
    $\mathrm{rev} \in L$ over one alphabet with at least two letters iff
    over every: $\mathrm{rev} \in L_\Sigma$ and a good $c : \Sigma
    \rightarrow \Gamma^\ell$, a good \emph{palindromic} $e : \Gamma
    \rightarrow \Sigma^k$ (Lemma~\ref{lem:palindrome}) and $q \triangleq c
    \circ e$ give
    \begin{align}
        \mathrm{rev}_\Gamma \;=\; H_q \circ \llbracket E' \rrbracket
        \circ q, \qquad E' = E^c .
    \end{align}
\end{theorem}

\begin{proof}
    $q(T) = c(e(T))$, so $\llbracket E' \rrbracket(q(T)) = c(\mathrm{rev}
    \, e(T))$ by Direction 1; the palindrome clause gives
    $\mathrm{rev}\, e(T) = e(\mathrm{rev}\, T)$, hence
    $\llbracket E' \rrbracket(q(T)) = c(e(\mathrm{rev}\,T)) = q(\mathrm{rev}
    \,T)$, and $H_q$ returns $\mathrm{rev}\,T$. The needed dictionaries
    exist for every pair of alphabets: from a source of at least three
    letters use the middle-marker family for $e$; from a binary source,
    first move to a ternary alphabet along the size-3 binary dictionary of
    Lemma~\ref{lem:palindrome}(iii) -- a coding of three letters needs
    three codewords, and three exist -- and then proceed as before; $c$
    itself is never required to be palindromic (Lemma~\ref{lem:dict}
    suffices), and $D_q = \{c(e(\gamma))\}$ is comma-free by
    Lemma~\ref{lem:compose}. (Verified: $e = \{\mathtt{x} \mapsto
    \mathtt{bab}, \mathtt{y} \mapsto \mathtt{cac}\}$ over $\Sigma = \{
    \mathtt{a},\mathtt{b},\mathtt{c}\}$, $c$ the length-6 double-$a$ coding
    to $\{\mathtt{x},\mathtt{y}\}$; $q$ and $H_q$ as \emph{expressions},
    with the oracle $D = c \circ \mathrm{rev} \circ c^{-1}$ standing in for
    $\llbracket E' \rrbracket$ ($\mathrm{rev}$-reachability being the
    hypothesis): $H_q(D(q(T))) = \mathrm{rev}\,T$ on all 63 inputs of
    length $\le 5$, zero failures.)
\end{proof}

\begin{remark}
    (\emph{No Contradiction with the Native Results})
    \label{rem:native}
    Theorem~\ref{thm:transfer} conjugates; Propositions
    \ref{prop:del-leftmost} and \ref{prop:kth} speak of \emph{native}
    functions, and the two live in different conjugacy classes. The binary
    witness $W$ of the leftmost deletion transfers to a ternary alphabet as
    the \emph{block-level} function $[\,\epsilon/c(b)\,]_1$ on the image of
    $c$ -- it deletes an $\ell$-block, not a character, and is not the
    native ternary $[\,\epsilon/\gamma\,]_1$; the negative ternary searches
    could not have seen the transferred witness in any case, its depth
    after transfer and retraction being far beyond the exhaustive depth-3
    and randomized $4$--$7$-pass spaces. Nor is the $\ell = 1$ coding a
    loop-hole:
    an injective character map is (trivially) good, and conjugating the
    binary once-deletion along it yields the native ternary function
    \emph{restricted to a two-letter sub-alphabet} -- reachable as a partial
    function, exactly the anchored-domain face of the phenomenon, and
    consistent with the resistance being about general position over three
    letters. Finally, the once-calculus -- and the rightmost and
    $k$-th-occurrence calculi -- have transfer theorems of their own:
    Lemma~\ref{lem:pass-transfer} is scan-direction- and
    multiplicity-agnostic, the first coded occurrence being the coding of
    the first native one. The marker-freshness obstruction of
    Proposition~\ref{prop:kth}(iii) is a statement about native simulation
    within one alphabet and is untouched.
\end{remark}

\subsection{The Unary Edge}

\begin{proposition}
    (\emph{The Boundary $|\Sigma| = 1$})
    \label{prop:unary-edge}
    \begin{enumerate}[(i)]
        \item $c(\mathtt{a}) = \mathtt{ab}$ is good ($\{w\}$ is
        comma-free iff $w$ is primitive), and every unary-reachable
        function is the restriction of a binary-reachable one.
        \item There is no injective monoid homomorphism $\{\mathtt{a},
        \mathtt{b}\}^* \rightarrow \{\mathtt{a}\}^*$: any $h$ has
        $h(\mathtt{ab}) = h(\mathtt{a})h(\mathtt{b}) =
        h(\mathtt{b})h(\mathtt{a}) = h(\mathtt{ba})$ while $\mathtt{ab} \ne
        \mathtt{ba}$. So no coding-conjugation crosses the $|\Sigma| = 1$
        boundary from above, and the transfer theorem is sharp: it holds
        for exactly the alphabets of at least two letters.
        \item Every unary $L$-expression computes a length map in the
        closure of $\{n \mapsto n,\ n \mapsto k\}$ under pointwise sum
        (concatenation) and $(g,r,p) \mapsto \big(n \mapsto r(n)
        \lfloor g(n)/p(n) \rfloor + g(n) \bmod p(n)\big)$ (the pass; the
        constant case is $n \mapsto i\lfloor n/j \rfloor + n \bmod j$, and
        $[\,X_1/a\,]X_1$ squares). Every such map is eventually
        polynomially bounded, so $n \mapsto 2^n$ is unreachable at any
        size. The floor-halving $n \mapsto \lfloor n/2 \rfloor$ -- the
        unary-domain function of Proposition~\ref{prop:unary-once} -- is
        absent from every unary expression of size $\le 7$ (36,978
        expressions, 2,091 distinct length maps on $n \le 24$): every
        single pass leaks its residue $n \bmod j$, and subtraction is not
        among the closures.
    \end{enumerate}
\end{proposition}

\begin{proof}
    (i) is Direction 1 for the coding $c$ (verified: all 258 depth-$\le 1$
    unary expressions on $a^0..a^{24}$, 6,450 evaluations, 0 mismatches,
    plus 12,000 deeper-random; alignment and the pass law, 566 cases). (ii)
    is the displayed computation. (iii): the pass law is the greedy scan of
    a unary string (verified: 20 $(i,j)$ pairs); the growth bound is an
    induction -- a pass outputs at most $|R| \cdot \lfloor |E|/|P| \rfloor
    \cdot |P| + |E| \le |R||E| + |E|$ characters; the halving search is
    exhaustive at size $\le 7$ ($\lceil n/2 \rceil$ and $n \bmod 2$ are
    present at size 4, $n^4$ at size 10 via $[\,[X/a]X / a\,]([X/a]X)$,
    verified). Of 480 total binary length profiles $n \mapsto
    |\llbracket E \rrbracket(a^n)|$ ($n \le 12$), exactly two are realized
    by no size-$\le 7$ unary expression nor any constant shift of one --
    bounded-search evidence for the strictness of (i)'s converse.
\end{proof}
